% Source: sources/returns/S6_04_MAPS_MANIFOLDS/sections/fourfolds_D2_c2.tex:1–336
\paragraph{Theorem D.2.35 (Proved: the fifteen minimal $c=2$ patterns).}
Take one of the fifteen patterns with $c=2$ in
\path{records/root_remaining_smooth_profiles.json}, use the primitive
Mumford filling at every multiplicative site, and make no additional
regular logarithmic transformations. Retain the additive original/log
choice of that pattern and the ordinary clutching $e\gamma$ of
(D.HY3). Every simply connected output, for every allowed finite
multiplicity and affine twist, is orientation-preservingly diffeomorphic
to the conic surgery $M_{\mathrm{conic}}$ of D.2.8. Equivalently, if
$\alpha\ne0$ in $H^2(X;\mathbb F_2)$, then
\[
 H_2(X;\mathbb Z)=H_3(X;\mathbb Z)=\mathbb Z/2,
 \qquad \langle\alpha^3,[X]_2\rangle=1.                 \tag{D.C21}
\]
The marked patterns, with zero positions numbered from zero as in the
source tables, are
\[
\begin{array}{c|r|r|l|l}
\text{OS/profile}&a&b&\text{zero positions}&\text{additive choice}\\\hline
47&7&2&2,3,4&\text{none}\\
47/II&7&2&2&II\text{ original or log}\\
49&1&6&2,3&IV^*\text{ log}\\
56&5&2&1&\text{none}\\
56&15&2&3,4&\text{none}\\
56/II&5&2&3&II\text{ original or log}\\
56/IV&1&6&0&IV\text{ log}\\
56/IV&5&6&3&IV\text{ log}\\
56/IV&15&2&3&IV\text{ original}.
\end{array}                                                   \tag{D.C22}
\]
Thus these fifteen patterns occupy thirteen distinct marked initial
rows. The statement includes the original additive caps in this table;
it does not replace the stipulated multiplicative Mumford caps by
other caps, or include additional regular logarithmic transformations.

\begin{proof}
\emph{1. The common piece and its integral markings.}
Throughout, use the actual lattice basis $(u,w,\delta,\gamma)$ and
write
\[
 \Omega=a\,u^*\wedge w^*+b\,\gamma^*\wedge\delta^*.
                                                               \tag{D.C23}
\]
Here $a$ is odd, $b=2$ or $6$, and $\gcd(a,b)=1$.
Remove the additive caps and one ordinary disc carrying the clutching;
fill every multiplicative site by its stated Mumford cap, and denote
the resulting common piece by $W$. Denote its regular torus-bundle
part by $V$. In particular all finite affine twists are outside $W$,
and $V$ has the linear section used in the marked construction.

There is an exact integral fact about these thirteen rows. The columns
of their actual multiplicative angular vanishing lattices, without
saturation, generate precisely
\[
 \Lambda_0=\ker\gamma^*=\mathbb Zu\oplus\mathbb Zw
                    \oplus\mathbb Z\delta.                    \tag{D.C24}
\]
The same is true of all columns of $A_i-I$. Both assertions are checked
by integer Hermite forms, not ranks or component labels, in the
certificate below. Filling a Mumford cap kills its base meridian and
its angular lattice. The usual torus-bundle presentation therefore
gives
\[
 \pi_1(W)=\mathbb Z\gamma\times F_{r-1},
                                                               \tag{D.C25}
\]
where $r\ge1$ is the number of unfilled discs. Every remaining
monodromy is the identity on the quotient $\mathbb Z\gamma$.
The section makes the displayed extension split. In this table
$r\le2$, but only torsion-freeness is needed below.

We also need an integral cohomology comparison. Each primitive
Mumford cap has a surjective degree-two boundary map, as proved in
D.2.26. If there are $N$ original punctures, then $b_1(V)=N$ by the
second assertion in (D.C24). A filled boundary has first Betti number
one greater than its cap. In the degree-one Mayer--Vietoris map the
target rank minus the source rank is consequently $r$, exactly
$b_1(W)$ by (D.C25). Its kernel has rank zero and is a subgroup of a
free group, hence is zero. Eliminating the surjective cap degree-two
maps now shows that
\[
 H_2(V;\mathbb Z)\longrightarrow H_2(W;\mathbb Z)
                  \text{ is onto}.                            \tag{D.C26}
\]
Both $H_1(V)$ and $H_1(W)$ are free. Thus restriction from $W$ to $V$
is injective in degree-two integral cohomology and also with
$\mathbb F_2$ coefficients. These facts will prevent a swept-class
ambiguity in the auxiliary-line calculation.

\emph{2. The actual symmetric line and its Chern class.}
On a multiplicative chart let $h_1$ be the elliptic period, let $h_G$
represent the section $G$, and let $R$ be the contraction multiplier.
The section pair is $P=aG,Q=bG$. On an auxiliary line with coordinate
$\zeta$ put
\[
\begin{aligned}
 h(z,v,\zeta)&=(h_1z,h_G^av,(-1)^a z^{-a}\zeta),\\
 F(z,v,\zeta)&=(h_G^{-b}z,Rv,
                 h_G^{-ab/2}R^{b/2}v^b\zeta),\\
 J(z,v,\zeta)&=(z^{-1},v^{-1},z^{-a}\zeta).
\end{aligned}                                                   \tag{D.C27}
\]
All exponents are integers. Direct multiplication gives
$hF=Fh$, $J^2=1$, $JhJ=h^{-1}$ and $JFJ=F^{-1}$; the integral
logarithmic constant $a$ in the third identity exponentiates to one.
The displayed $J$ is normalized on the rigidifying origin.
The alternating Chern form is $\Omega$. The normalized half-period
signs are $(-1)^a$ on $u,w$ and $+1$ on $\delta,\gamma$; thus its
mod-two semicharacter is
\[
                    q=a(u+w+uw).                              \tag{D.C28}
\]
Every actual monodromy preserves this semicharacter. The certificate
checks this on all sixteen vectors modulo two at every marked site.
The symmetric rigidified Appell--Humbert identifications consequently
patch on overlap markings: an isomorphism preserving the rigidification
is unique, so their compositions satisfy the cocycle condition.
This constructs the holomorphic auxiliary line on $V$ from the actual
factors (D.C27).
Since $b$ is even, the half-$\delta$ translation lifts in these factors
as $(z,v,\zeta)\mapsto(z,-v,\zeta)$. It commutes with $h,F$.
The overlap identifications preserve this lift: in the invariant
prequantum marking the lift depends only on $\Omega$ and $\delta/2$,
and each monodromy fixes $\delta$ and preserves $\Omega$ and $q$.
Thus the line has the global involution linearization used below.

Compare it with the flat prequantum family having connection potential
$\pi i\Omega(x,dx)$ and the same semicharacter, as in the fully marked
comparison in D.2.8. The monodromies preserve this family. Both lines
are symmetric and rigidified with alternating form $\Omega$; their
ratio is therefore fibrewise flat of order dividing two. The normalized
square-trivializations on the fibres are unique and patch, making the
ratio globally two-torsion. Since $H^2(V;\mathbb Z)$ is free, its Chern
class is zero. The actual line has precisely the pure fibre class
$\Omega$ and zero value on the marked swept cycles. This argument
controls the mixed values as well as the alternating fibre form.

These factors extend to a line $\mathcal L$ on $W$. Here is the
integral local extension, including the two parities that we need.
In a genuine $SL_2(\mathbb Z)$ elliptic marking the actual Tate
valuations are
\[
 (\operatorname{ord}h_1,\operatorname{ord}h_G,
       \operatorname{ord}R)=(n,j,\rho),\qquad
 P_{\rm val}=\begin{pmatrix}n&-bj\\aj&\rho\end{pmatrix}.
                                                               \tag{D.C29}
\]
The determinant is one at the zero and zero at every other
multiplicative site. These are raw valuations from the actual
monodromy; no residue of $j$ modulo $n$ is substituted.
At the zero, assign the following integer values to the lattice
vertices and interpolate linearly on the given unimodular fan:
\[
 \varphi(x,y)=
 \frac{-a\rho x^2-2abjxy+bn y^2+ax}{2}.                        \tag{D.C210}
\]
Indeed $n\rho+abj^2=1$, so $n,\rho$ are odd, and the expression is
integral at every integer vertex. Its two period differences are
\[
\begin{aligned}
 \varphi(x+n,y+aj)-\varphi(x,y)&=-ax,\\
 \varphi(x-bj,y+\rho)-\varphi(x,y)
          &=by+(b\rho-abj)/2.
\end{aligned}                                                   \tag{D.C211}
\]
They are exactly the valuations of the line factors in (D.C27).
The usual monomial frames with these integral slopes therefore extend
the line and its commuting lattice identifications over every toric
chart. This uses an integral piecewise affine support function; no
positivity of the auxiliary line is required.

At a rank-one point write the two valuation periods as
$H=\alpha\nu$, $F=\beta\nu$, where
$\nu=(\nu_1,\nu_2)$ is primitive and
$\gcd(\alpha,\beta)=1$. Thus
$n=\alpha\nu_1$, $aj=\alpha\nu_2$,
$-bj=\beta\nu_1$ and $\rho=\beta\nu_2$.
The identity $-a\beta\nu_1=b\alpha\nu_2$ shows that
$k=a\nu_1/\alpha$ is an integer. On the valuation line use
\[
             \varphi_1(t)=-k t(t-\alpha)/2.
\]
Its two exact period differences are
$-a\nu_1t$ and
$b\nu_2t+(b\rho-abj)/2$, respectively, as required by
(D.C27). At every odd-isotropy site $\nu_1$ is odd and
$\beta\nu_1=-bj$ is even, so $\beta$ is even and $\alpha$
is odd. The values of $\varphi_1$ at integer vertices are therefore
integral. At the unique even-isotropy site the valuation vector is
$(2,aj)$ with $aj$ odd, so $\alpha=1$ and the same integrality
holds. Integral monomial frames on the one-dimensional fan, retaining
the invertible second short coordinate, give the required extension.
The two displayed period identities prove that both actual deck
factors descend this extension; it is not enough to match only one
of them.
Holomorphic unit factors extend in the same frames and affect neither
integral slope nor stabilizer character. The local extensions agree
with the line on their annuli and hence glue to $\mathcal L$ on $W$.
By (D.C26), its Chern class has exactly the pure fibre and zero swept
values already computed on $V$.

\emph{3. The fixed curves and the two line characters.}
Let $\iota$ be the scalar involution $(z,v)\mapsto(z,-v)$.
The exact primitive angular vectors show that its fixed locus on $W$
is the disjoint union of one elliptic curve $E$ at the unique $I_2$
site and the rational vertical-edge curve $C$ at the rank-two zero.
All other rank-one scalar stabilizers have odd order; this is checked
from the actual vanishing vectors in the certificate. The zero has
valuation lattice all of $\mathbb Z^2$, and $\delta=(0,1)$ in that
short lattice. Thus there is one vertical-edge orbit, with the same
normal model
$\mathcal O(-1)_{+1}\oplus\mathcal O(-1)_{-1}$ as in D.2.14.

At the zero the vertical slope difference is
\[
 \varphi(x,y+1)-\varphi(x,y)=-abjx+bny+bn/2.
\]
Its two nonconstant coefficients are even and its constant is odd,
since $b/2$ and $n$ are odd. The induced line character along $C$
is therefore nontrivial. The degree along its adjacent toric wall is,
up to the immaterial orientation convention for the divisor,
$-a\rho+abj$, which is odd. The other diagonal convention changes
this expression by an even integer, so
\[
 \operatorname{char}_{\iota}(\mathcal L|_C)=1,
          \qquad \deg(\mathcal L|_C)\equiv1\pmod2.            \tag{D.C212}
\]
At the $I_2$ site, $j$ is odd and $\rho=-abj^2/2$. The coordinates
\[
 x=z^{(1-aj)/2}v,\qquad y=z^{-aj}v^2
\]
are an integral unimodular change of the short torus coordinates;
they send the primitive vanishing vector $(2,aj)$ to $(1,0)$.
The scalar involution acts as $(-x,y)$, giving sign $-1$ on both
normal directions to $E$. The $h$ multiplier is, up to a unit,
$(-1)^a x^{-2a}y^a$. After replacing $F$ by $Fh^{bj/2}$, its
$x$ exponent is zero and its $y$ exponent is $b/2$. The support
function $-at(t-1)$ has all slopes even. Hence
\[
                 \operatorname{char}_{\iota}(\mathcal L|_E)=0.
                                                               \tag{D.C213}
\]

There are no further fixed points in the removed caps when they are
reattached. An original additive cap of this table has zero line
multidegree and a free scalar circle, by its actual product description
in D.2.26. On a finite cap, a scalar fixed-point equation on the covering
torus would imply $jn_i/m_i=0$ in $\mathbb R/\mathbb Z$ on applying
$\gamma^*$. Simple connectivity requires $\gcd(m_i,n_i)=1$, so the
root power is trivial; translation by $\delta/2$ cannot then fix a
torus point. Away from the central torus the base rotation gives the
same conclusion. The ordinary clutch disc also has free scalar action.
Thus $X^{\iota}=C\sqcup E$ for every multiplicity and twist in the
statement, including even finite multiplicities.

\emph{4. The global mod-two attachment.}
D.2.26 gives $H_2(X)=H_3(X)=\mathbb Z/2$, generated in degree two by
the image of the primitive fibre class $s$; there are no additional
regular sweeps. We verify which cohomology class this is on $W$.
For every finite log in this table, $b/q_i=2$. Its complete two normal
columns (D.HY7) have pure fibre coefficients
\[
                 bn_i/q_i=2n_i,
                \qquad -\alpha_i b r_i,                       \tag{D.C214}
\]
where $\alpha_i m_i+\beta_i n_i=1$. Both coefficients are even,
for arbitrary allowed $m_i,n_i,r_i$. The original additive cap normal
columns have zero fibre coefficient in the rigidified marking.
The ordinary clutching $e\gamma$ has fibre coefficients divisible by
$b$, and its global relation is $be\,s$. Thus the functional that is
one on $s$ and zero on all swept generators annihilates every normal
column modulo two. It induces the nonzero class $\alpha$ on $X$.
This uses the full integral normal columns, not a transfer of even
degree. Its restriction to $V$ is exactly the reduction of the Chern
class just constructed, and (D.C26) now implies
\[
                  \alpha|_W=\rho_2c_1(\mathcal L).             \tag{D.C215}
\]
The absence of an $\operatorname{Ext}(H_1(W),\mathbb F_2)$ term is
part of (D.C25), not an implicit universal-coefficient choice.

\emph{5. The full two-primary smooth invariant.}
Since $H^2(X;\mathbb Z)=0$, the complex manifold has $c_1=0$ and is
spin. Put $G=\langle\iota\rangle$ and let
$\tau\in H^1(BG;\mathbb F_2)$ be its generator. As in D.2.8,
$\alpha$ has a Borel equivariant lift $\widetilde\alpha$:
$H^1(X;\mathbb F_2)=0$ kills the first possible obstruction, and the
fixed point splits the map to $BG$, killing the obstruction into
$H^3(BG;\mathbb F_2)$.

We check the relative fixed-point normalization on $W$. The scalar
circle acts trivially on its fundamental group, by an isotopy fixing
a point of $C$. That point splits the Borel extension, so
\[
                 \pi_1(W_G)=\pi_1(W)\times G.
\]
The first factor is torsion-free by (D.C25). Consequently every fixed
point gives the same order-two element $(1,\iota)$ in this product.
By (D.C215), the difference between $\widetilde\alpha|_{W_G}$ and
the equivariant Chern class of $\mathcal L$ vanishes in ordinary
cohomology. Pull it to the universal cover and use the degree-two
Cartan--Leray sequence: the difference comes from the group cohomology
of $\pi_1(W)\times G$. Its possible contributions are a common
multiple of $\tau^2$ and a mixed class $\tau\xi$ with
$\xi\in H^1(W;\mathbb F_2)$; its ordinary group contribution is
zero, since $H^2(B\pi_1(W))\to H^2(W)$ is injective. After one global
adjustment by $\tau^2$, (D.C212)--(D.C213) give
\[
 \widetilde\alpha|_C=\alpha|_C+\tau^2,
 \qquad \widetilde\alpha|_E=\alpha|_E+\tau\eta,
       \quad\eta\in H^1(E;\mathbb F_2).                       \tag{D.C216}
\]
In particular two unrelated fixed-point normalizations have not been
chosen.

Both complex normal directions have sign $-1$. Their normal first
Chern class is even because $c_1(X)=0$ and the tangent Chern class of
each fixed curve is even. Their equivariant Euler class is therefore
$\tau^4$. Equivariant Thom localization yields
\[
 \langle\alpha^3,[X]_2\rangle
 =\int_C\frac{(\alpha|_C+\tau^2)^3}{\tau^4}
  +\int_E\frac{(\alpha|_E+\tau\eta)^3}{\tau^4}
 =\int_C\alpha|_C=1.
\]
The elliptic contribution is zero since degree-one classes on an
oriented elliptic curve square to zero. Finally the exact primary
calculation
$\Omega^{\mathrm{Spin}}_6(K(\mathbb Z/2,2))=\mathbb Z/2$,
detected by this cube, and Zhubr's complete classification identify
$X$ with $M_{\mathrm{conic}}$; the classification and the conic's cube
were established in D.2.8 directly from
\cite[Theorems 3, 2.10 and Section 3.12]{manZhubr75}.
This retains the two-primary refined invariant, not just ordinary
Pontryagin classes.
\end{proof}
